% Motor: XeLaTeX, fijado en .latexmkrc y en la receta del editor.
\documentclass[aspectratio=169,10pt,c]{beamer}

% !TEX root = slides.tex
% Preambulo comun: estilo UTP, colores, fuentes, plantillas de beamer y
% comandos de apoyo didactico. Se carga desde slides.tex tras \documentclass.

% La opción "c" de beamer reparte el espacio libre 1:1.5 (arriba:abajo), por lo
% que el contenido queda por encima del centro. Igualamos ambos rellenos.
\makeatletter
\define@key{beamerframe}{c}[true]{%
  \beamer@frametopskip=0pt plus 1fill\relax%
  \beamer@framebottomskip=0pt plus 1fill\relax%
  \beamer@frametopskipautobreak=0pt plus .5\paperheight\relax%
  \beamer@framebottomskipautobreak=0pt plus .5\paperheight\relax%
}
\makeatother

\usepackage{graphicx}
\usepackage{tikz}
\usepackage{xcolor}
\usepackage{iftex}

\ifPDFTeX
  \usepackage[utf8]{inputenc}
  \usepackage[T1]{fontenc}
  \usepackage{lmodern}          % tipografia escalable para el respaldo con pdflatex
  \renewcommand{\familydefault}{\sfdefault}
\else
  \usepackage{fontspec}
  \IfFontExistsTF{Arial}
    {\setsansfont{Arial}[BoldFont={Arial Bold}]}
    {\setsansfont{Helvetica}[BoldFont={Helvetica Bold}]}
\fi

\usepackage[spanish,es-nodecimaldot]{babel}
\usepackage{microtype}
\usepackage{booktabs}
\usepackage{amssymb}
\usepackage{array}
\usetikzlibrary{arrows.meta,positioning,fit,calc,shapes.geometric}

\graphicspath{{assets/}}

% ---------- Identidad visual UTP ----------
\definecolor{utpRojo}{HTML}{C42B35}
\definecolor{utpNegro}{HTML}{111111}
\definecolor{utpGris}{HTML}{5F5F5F}
\definecolor{utpGrisMedio}{HTML}{9A9A9A}
\definecolor{utpGrisClaro}{HTML}{F1F1F1}
\definecolor{utpNaranja}{HTML}{C1670B}

\usetheme{default}
\setbeamercolor{normal text}{fg=utpNegro,bg=white}
\setbeamercolor{structure}{fg=utpRojo}
\setbeamercolor{frametitle}{fg=utpNegro,bg=white}
\setbeamercolor{title}{fg=utpNegro}
\setbeamercolor{subtitle}{fg=utpNegro}
\setbeamercolor{author}{fg=utpNegro}
\setbeamercolor{date}{fg=utpNegro}
\setbeamercolor{block title}{fg=white,bg=utpNegro}
\setbeamercolor{block body}{fg=utpNegro,bg=utpGrisClaro}
\setbeamercolor{alertblock title}{fg=white,bg=utpRojo}
\setbeamercolor{alertblock body}{fg=utpNegro,bg=utpRojo!7}
\setbeamercolor{exampleblock title}{fg=white,bg=utpGris}
\setbeamercolor{exampleblock body}{fg=utpNegro,bg=utpGris!8}

\setbeamerfont{title}{series=\bfseries,size=\fontsize{19.5}{22.5}\selectfont}
\setbeamerfont{subtitle}{size=\fontsize{11.5}{13.5}\selectfont}
\setbeamerfont{date}{series=\bfseries,size=\fontsize{11.5}{13.5}\selectfont}
\setbeamerfont{author}{size=\fontsize{10.5}{12.5}\selectfont}
\setbeamerfont{frametitle}{series=\bfseries,size=\large}
\setbeamerfont{block title}{series=\bfseries,size=\small}
\setbeamerfont{block body}{size=\normalsize}

\setbeamertemplate{navigation symbols}{}
\setbeamertemplate{headline}{}
\setbeamertemplate{blocks}[default]
\setbeamertemplate{itemize item}{\color{utpRojo}\small\raise.15ex\hbox{$\blacktriangleright$}}
\setbeamertemplate{itemize subitem}{\color{utpGris}\scriptsize$\bullet$}
\setbeamertemplate{enumerate item}{\color{utpRojo}\insertenumlabel.}
\setbeamertemplate{background canvas}{\color{white}\rule{\paperwidth}{\paperheight}}

% Título de diapositiva: negrita negra a la izquierda, logo UTP a la derecha
\setbeamertemplate{frametitle}{%
  \begin{tikzpicture}[remember picture,overlay]
    \node[anchor=north east,inner sep=0pt]
      at ([xshift=-7mm,yshift=-5.5mm]current page.north east)
      {\includegraphics[width=17mm]{utp-logo.png}};
  \end{tikzpicture}%
  \vspace{1.2mm}%
  \begin{minipage}[t]{.80\textwidth}\raggedright\insertframetitle\end{minipage}\par
  \vspace{1.1mm}{\color{utpRojo}\rule{\textwidth}{1.1pt}}
}

% Solo numeración discreta; la plantilla UTP no lleva pie de página
\setbeamertemplate{footline}{%
  \hfill{\color{utpGrisMedio}\scriptsize\insertframenumber}\hspace{6mm}\vspace{2.5mm}%
}

% Portada UTP
\setbeamertemplate{title page}{%
  \begin{tikzpicture}[remember picture,overlay]
    \node[anchor=center,inner sep=0pt] at (current page.center) {%
      \includegraphics[width=\paperwidth,height=\paperheight]{utp-cover-bg.jpeg}%
    };
    \node[anchor=north,align=center,text=utpNegro,inner sep=0pt,
      text width=145mm,font=\usebeamerfont{title}]
      at ([yshift=-25.2mm]current page.north) {\inserttitle};
    \node[anchor=north,align=center,text=utpNegro,inner sep=0pt,
      text width=142mm,font=\usebeamerfont{subtitle}]
      at ([yshift=-38.2mm]current page.north) {\insertsubtitle};
    \node[anchor=north,align=center,text=utpNegro,inner sep=0pt,
      font=\usebeamerfont{date}]
      at ([yshift=-50.4mm]current page.north) {\insertdate};
    \node[anchor=north,align=center,text=utpNegro,inner sep=0pt,
      font=\usebeamerfont{author}]
      at ([yshift=-58.4mm]current page.north) {\insertauthor};
  \end{tikzpicture}%
}

% ---------- Comandos de apoyo didáctico ----------
\newcommand{\concepto}[2]{%
  \begin{block}{#1}\large #2\end{block}}
\newcommand{\pregunta}[1]{%
  \begin{alertblock}{Decide y explica}#1\end{alertblock}}
\newcommand{\analogia}[1]{%
  \begin{exampleblock}{En palabras simples}\small #1\end{exampleblock}}
\newcommand{\enlarealidad}[1]{%
  \begin{block}{Cómo se ve en la realidad}\small #1\end{block}}
\newcommand{\ojo}[1]{%
  \begin{alertblock}{Ojo con esto}\small #1\end{alertblock}}
\newcommand{\apoyo}{{\normalfont\footnotesize\color{utpGris}~· material de apoyo}}
\newcommand{\marca}[1]{\tikz[baseline=(m.base)]{%
  \node[rounded corners=1.5pt,fill=utpGrisClaro,draw=utpGrisMedio,inner xsep=3pt,inner ysep=1.5pt,
        font=\footnotesize] (m) {#1};}}

\tikzset{
  caja/.style={draw=utpNegro, very thick, rounded corners=1.5pt, fill=white,
    align=center, minimum height=0.72cm, inner sep=5pt, font=\small},
  cajaazul/.style={caja, draw=utpGris, fill=white},
  cajaverde/.style={caja, draw=utpNegro, fill=utpGrisClaro},
  cajanaranja/.style={caja, draw=utpRojo!75, fill=utpRojo!7},
  flecha/.style={-{Stealth[length=2.3mm]}, very thick, draw=utpNegro},
  flechavuelta/.style={{Stealth[length=2.3mm]}-{Stealth[length=2.3mm]}, very thick, draw=utpNegro}
}

% ---------- Notación pata de gallo (crow's foot) ----------
% La rotación indica hacia dónde está la entidad: 0 derecha, 180 izquierda,
% 90 arriba, -90 abajo. La pata se abre hacia la entidad del lado "muchos".
\tikzset{
  relacion/.style={thick, draw=utpNegro},
  pics/uno/.style={code={%
    \draw[thick,utpNegro,line cap=round] (0,-1.5mm) -- (0,1.5mm);}},
  pics/muchos/.style={code={%
    \draw[thick,utpNegro,line cap=round,line join=round]
      (2.6mm,1.5mm) -- (0,0) -- (2.6mm,-1.5mm);
    \draw[thick,utpNegro,line cap=round] (0,0) -- (2.6mm,0);}},
}
\newcommand{\simbuno}{\tikz[baseline=-0.6ex]{\draw[relacion](0,0)--(7mm,0);\pic at (4.4mm,0){uno};}}
\newcommand{\simbmuchos}{\tikz[baseline=-0.6ex]{\draw[relacion](0,0)--(7mm,0);\pic at (4.4mm,0){muchos};}}


\title{Claves y restricciones en SQL}
\subtitle{Del modelo E-R al esquema relacional en DDL:\\ claves primarias, restricciones y claves foráneas}
\date{Semana 06 - Sesión 12}
\author{Rel Guzman}

% Espacio en blanco para completar
\newcommand{\bl}[1][1.9cm]{\underline{\hspace{#1}}}
% Nombre de tabla, siempre en naranja
\newcommand{\tbl}[1]{\textcolor{utpNaranja}{\texttt{#1}}}
% Palabra reservada de SQL dentro de un bloque de código
\newcommand{\kw}[1]{\textcolor{utpRojo}{\textbf{#1}}}
% Parte nueva de un bloque de codigo que crece paso a paso
\newcommand{\nuevo}[1]{{%
  \setlength{\fboxsep}{1.2pt}%
  \begingroup
  \renewcommand{\kw}[1]{\textcolor{utpNegro}{\textbf{##1}}}%
  \colorbox{yellow!45}{\textcolor{utpNegro}{#1}}%
  \endgroup}}
% Marcas de clave
\newcommand{\pk}{{\scriptsize\color{utpRojo}\textbf{PK}}}
\newcommand{\fk}{{\scriptsize\color{utpNaranja}\textbf{FK}}}
% Insignia de práctica
\newcommand{\ins}[1]{\tikz[baseline=-0.5ex]{\node[rounded corners=2pt,fill=utpRojo,text=white,
  font=\bfseries\scriptsize,inner xsep=5pt,inner ysep=2.5pt] {#1};}}
% Bloque de código SQL sobre fondo gris
\newsavebox{\sqlbuf}
\newenvironment{sqlcode}[1][.86\textwidth]
  {\begin{center}\begin{lrbox}{\sqlbuf}\begin{minipage}{#1}\ttfamily\footnotesize\vspace{6pt}\raggedright}
  {\vspace{6pt}\end{minipage}\end{lrbox}\colorbox{utpGrisClaro}{\usebox{\sqlbuf}}\end{center}}
% Sangría dentro de un bloque de código
\newcommand{\ind}{\hspace*{1.6em}}
% Tabla de tres columnas para alinear nombre / tipo / restricciones
\newenvironment{sqltab}
  {\begin{tabular}{@{}l@{\hspace{1.4em}}l@{\hspace{1.4em}}l@{}}}
  {\end{tabular}}
\newcommand{\sqlline}[1]{\multicolumn{3}{@{}l}{#1}}

\begin{document}

% 1 ---------------------------------------------------------------------
{\setbeamertemplate{footline}{}\maketitle}

% 2 ---------------------------------------------------------------------
\begin{frame}[c]
\begin{block}{Logro de la sesión}
Al finalizar, podrás definir \textbf{claves primarias} y \textbf{restricciones}, crear una base de datos y conectar sus tablas mediante \textbf{claves foráneas}.
\end{block}
\end{frame}

% 3 ---------------------------------------------------------------------
\begin{frame}{Dónde estamos}
\begin{center}
\resizebox{\textwidth}{!}{%
\begin{tikzpicture}[font=\small]
  \node[cajaverde,text width=3.0cm] (a) {Caso de estudio};
  \node[caja,text width=3.0cm,right=12mm of a] (b) {Diseño conceptual};
  \node[caja,text width=3.0cm,right=12mm of b] (c) {Diseño lógico};
  \node[cajanaranja,text width=3.0cm,right=12mm of c] (d)
    {\textbf{Diseño físico}\\[2pt] \texttt{CREATE TABLE}};
  \draw[flecha] (a)--(b);
  \draw[flecha] (b)--(c);
  \draw[flecha] (c)--(d);
  \node[font=\tiny,below=2mm of b,text=utpGris] {entidades y relaciones};
  \node[font=\tiny,below=2mm of c,text=utpGris] {tablas, \pk{} y \fk{}};
\end{tikzpicture}}
\end{center}
\vspace{0.75cm}
\begin{block}{DDL}
\small DDL es el \textbf{lenguaje de definición de datos}: permite crear y modificar la estructura de la base de datos. \texttt{CREATE TABLE} crea sus tablas.
\end{block}
\end{frame}

% 4 ---------------------------------------------------------------------
\begin{frame}[plain]
  \begin{tikzpicture}[remember picture,overlay]
    \node[anchor=north east,inner sep=0pt]
      at ([xshift=-10mm,yshift=-8mm]current page.north east)
      {\includegraphics[width=43mm]{utp-logo.png}};
    \node[anchor=center,inner sep=0pt] at (current page.center)
      {\resizebox{100mm}{!}{\bfseries\color{utpNegro}Claves y restricciones}};
    \node[anchor=north,inner sep=0pt] at ([yshift=-11mm]current page.center)
      {\color{utpRojo}\rule{100mm}{1.6pt}};
  \end{tikzpicture}
\end{frame}


\begin{frame}{Paso 4: ¿quién es cada fila?}
\begin{columns}[c,onlytextwidth]
\column{.58\textwidth}
\begin{sqlcode}[\textwidth]
\kw{CREATE TABLE} cliente\\
(\\
\ind \nuevo{id \ \ \ \ \ INT \ \ \ \ \ \ \ \ \kw{PRIMARY KEY},}\\
\ind nombre \ VARCHAR(80),\\
\ind correo \ VARCHAR(120)\\
);
\end{sqlcode}
\column{.38\textwidth}
\begin{center}
\begin{tikzpicture}
  \node[cajaverde,text width=2.7cm,font=\scriptsize,align=left]
    {\textbf{cliente}\\[3pt] \texttt{id} \ \pk\\ \texttt{nombre}\\ \texttt{correo}};
\end{tikzpicture}
\end{center}
\vspace{5pt}
\centering\scriptsize la \pk{} del dibujo\\ es esta línea
\end{columns}
\vspace{0.5cm}
\small Sin el \texttt{id} hay un problema: si dos clientes se llaman igual, \textbf{no se pueden distinguir}. La \textbf{clave primaria} es la columna que identifica a una sola fila; no se repite y nunca queda vacía.
\vspace{0.2cm}

\small Y con eso la instrucción está completa: \textbf{eso es todo un \texttt{CREATE TABLE}}.
\end{frame}

% 13 ---------------------------------------------------------------------
\begin{frame}{La instrucción completa, con sus nombres}
\begin{sqlcode}[.60\textwidth]
\begin{sqltab}
\sqlline{\kw{CREATE TABLE} cliente (}\\
\ind id     & INT           & \kw{PRIMARY KEY},\\
\ind nombre & VARCHAR(80)   & \kw{NOT NULL},\\
\ind correo & VARCHAR(120)  & \kw{UNIQUE}\\
\sqlline{);}\\
\end{sqltab}
\end{sqlcode}
\vspace{0.25cm}
\begin{center}
\small
\begin{tabular}{>{\ttfamily}p{3.1cm}>{\raggedright\arraybackslash}p{8.0cm}}
\toprule
\normalfont\bfseries Parte & \bfseries Cómo se llama \\
\midrule
cliente & el \textbf{nombre de la tabla} \\
nombre & el \textbf{nombre de la columna} \\
VARCHAR(80) & el \textbf{tipo de dato} \\
PRIMARY KEY & la \textbf{clave primaria} \\
NOT NULL, UNIQUE & las \textbf{restricciones} \\
\bottomrule
\end{tabular}
\end{center}
\vspace{0.25cm}
\small Faltan las dos últimas, y son las que siguen.
\end{frame}

% 14 ---------------------------------------------------------------------
\begin{frame}{El tipo dice qué clase de cosa se guarda}
\begin{center}
\begin{tikzpicture}
  \node[caja,text width=2.4cm,font=\scriptsize,align=center] (a) {89.00};
  \node[caja,text width=2.4cm,font=\scriptsize,align=center,right=10mm of a] (b) {Raúl};
  \node[caja,text width=2.4cm,font=\scriptsize,align=center,right=10mm of b] (c) {2026-03-04};
  \node[cajanaranja,text width=2.4cm,font=\scriptsize,align=center,below=9mm of a] (a2) {\texttt{DECIMAL(10,2)}};
  \node[cajanaranja,text width=2.4cm,font=\scriptsize,align=center,below=9mm of b] (b2) {\texttt{VARCHAR(80)}};
  \node[cajanaranja,text width=2.4cm,font=\scriptsize,align=center,below=9mm of c] (c2) {\texttt{DATE}};
  \draw[flecha] (a)--(a2); \draw[flecha] (b)--(b2); \draw[flecha] (c)--(c2);
\end{tikzpicture}
\end{center}
\vspace{0.3cm}
\begin{center}
\footnotesize
\begin{tabular}{>{\ttfamily}p{2.7cm}>{\raggedright\arraybackslash}p{4.4cm}>{\raggedright\arraybackslash}p{4.2cm}}
\toprule
\normalfont\bfseries Tipo & \bfseries Para qué & \bfseries Ejemplo \\
\midrule
INT & Números enteros & \texttt{id}, \texttt{cantidad} \\
VARCHAR(n) & Texto, hasta \texttt{n} caracteres & \texttt{nombre VARCHAR(80)} \\
DECIMAL(p,d) & Dinero & \texttt{precio DECIMAL(10,2)} \\
DATE & Fecha & \texttt{fecha DATE} \\
BOOLEAN & Verdadero o falso & \texttt{activo BOOLEAN} \\
\bottomrule
\end{tabular}
\end{center}
\end{frame}

% 15 ---------------------------------------------------------------------
\begin{frame}{\ins{Práctica 2} Elige el tipo de cada columna}
\begin{center}
\small
\begin{tabular}{>{\raggedright\arraybackslash}p{5.4cm}>{\raggedright\arraybackslash}p{3.0cm}>{\raggedright\arraybackslash}p{3.4cm}}
\toprule
\bfseries Dato que hay que guardar & \bfseries Tipo & \bfseries ¿Por qué? \\
\midrule
El precio de un producto & \bl[2.4cm] & \bl[2.8cm] \\[6pt]
El DNI del cliente & \bl[2.4cm] & \bl[2.8cm] \\[6pt]
La fecha del pedido & \bl[2.4cm] & \bl[2.8cm] \\[6pt]
Cuántas unidades se pidieron & \bl[2.4cm] & \bl[2.8cm] \\[6pt]
Si el producto sigue a la venta & \bl[2.4cm] & \bl[2.8cm] \\
\bottomrule
\end{tabular}
\end{center}
\vspace{0.35cm}
\small \textbf{Pista:} pregúntate dos cosas. ¿Se harán \textbf{cuentas} con ese dato? ¿Puede empezar con \textbf{cero}?
\end{frame}

% 16 ---------------------------------------------------------------------
\begin{frame}{Las restricciones: reglas que toda fila cumple}
\begin{center}
\small
\begin{tabular}{>{\ttfamily}p{2.9cm}>{\raggedright\arraybackslash}p{4.2cm}>{\raggedright\arraybackslash}p{4.6cm}}
\toprule
\normalfont\bfseries Restricción & \bfseries Qué obliga & \bfseries En la tienda \\
\midrule
NOT NULL & No puede quedar vacía & Un cliente sin nombre no sirve \\[5pt]
UNIQUE & No se repite en dos filas & Dos clientes no comparten correo \\[5pt]
AUTO\_INCREMENT & El gestor numera solo & \texttt{id} 1, 2, 3\ldots sin pensarlo \\[5pt]
DEFAULT \textit{valor} & Si no se indica, usa ese & \texttt{activo BOOLEAN DEFAULT TRUE} \\
\bottomrule
\end{tabular}
\end{center}
\vspace{0.4cm}
\small Se escriben \textbf{después del tipo}, en la misma línea. Son opcionales: una columna con nombre y tipo ya funciona.
\vspace{0.3cm}

\small \texttt{PRIMARY KEY} ya incluye \texttt{NOT NULL} y \texttt{UNIQUE}; por eso casi siempre se escribe junto con \texttt{AUTO\_INCREMENT} y nada más.
\end{frame}

% 17 ---------------------------------------------------------------------
\begin{frame}[plain]
  \begin{tikzpicture}[remember picture,overlay]
    \node[anchor=north west,font=\bfseries\fontsize{24}{28}\selectfont]
      at ([xshift=12mm,yshift=-9mm]current page.north west) {Preguntas};
    \node[anchor=north east,inner sep=0pt]
      at ([xshift=-10mm,yshift=-8mm]current page.north east)
      {\includegraphics[width=43mm]{utp-logo.png}};
    \node[anchor=center,inner sep=0pt]
      at ([yshift=-5mm]current page.center)
      {\includegraphics[width=78mm]{preguntas-ilustracion.png}};
  \end{tikzpicture}
\end{frame}

% 18 ---------------------------------------------------------------------
\begin{frame}{Antes de las tablas: crear la base y entrar en ella}
\begin{center}
\begin{tikzpicture}
  \node[caja,text width=7.0cm,minimum height=1.9cm,font=\scriptsize,align=center,
        label={[font=\scriptsize\bfseries,anchor=north west,xshift=3mm,yshift=-2mm]north west:tienda}] (bd) {};
  \node[cajaverde,text width=1.3cm,font=\tiny,align=center,xshift=-22mm,yshift=-1mm] at (bd) {cliente};
  \node[cajaverde,text width=1.3cm,font=\tiny,align=center,yshift=-1mm] at (bd) {pedido};
  \node[cajaverde,text width=1.3cm,font=\tiny,align=center,xshift=22mm,yshift=-1mm] at (bd) {producto};
  \node[font=\scriptsize,below=2mm of bd] {la base de datos es la \textbf{caja}; las tablas van dentro};
\end{tikzpicture}
\end{center}
\vspace{0.3cm}
\begin{sqlcode}[.72\textwidth]
\kw{CREATE DATABASE} tienda;\\
\kw{USE} tienda;\\[5pt]
\kw{CREATE TABLE} cliente ( \ldots );
\end{sqlcode}
\vspace{0.05cm}
\small \texttt{USE} dice en cuál se trabaja: desde esa línea, todo lo que se escriba se crea ahí dentro. Si se olvida, el gestor responde \texttt{No database selected}.
\vspace{0.2cm}

\analogia{Es el «abrir carpeta» antes de guardar. Si no abres ninguna, el archivo aparece en cualquier parte.}
\end{frame}

% 19 ---------------------------------------------------------------------
\begin{frame}{\texttt{FOREIGN KEY}: la conexión, por fin escrita}
\begin{center}
\begin{tikzpicture}
  \node[cajaverde,text width=1.9cm,font=\tiny,align=center] (cli) {\textbf{cliente}\\ \texttt{id} \pk};
  \node[cajanaranja,text width=2.1cm,font=\tiny,align=center,right=22mm of cli] (ped) {\textbf{pedido}\\ \texttt{cliente\_id} \fk};
  \draw[relacion] (cli)--(ped);
  \pic at ($(cli.east)!2.6mm!(ped.west)$) {uno};
  \pic at ($(ped.west)!2.6mm!(cli.east)$) {muchos};
\end{tikzpicture}
\end{center}
\vspace{0.2cm}
\begin{sqlcode}[.72\textwidth]
\begin{sqltab}
\sqlline{\kw{CREATE TABLE} pedido (}\\
\ind id          & INT  & \kw{PRIMARY KEY} \kw{AUTO\_INCREMENT},\\
\ind fecha       & DATE & \kw{NOT NULL},\\
\ind cliente\_id  & INT  & \kw{NOT NULL},\\
\sqlline{\ind \nuevo{\kw{FOREIGN KEY} (cliente\_id) \kw{REFERENCES} cliente(id)}}\\
\sqlline{);}\\
\end{sqltab}
\end{sqlcode}
\vspace{0.15cm}
\small \textbf{Son dos líneas, no una:} arriba se declara \texttt{cliente\_id} como cualquier columna; abajo se dice \textbf{a dónde apunta}. Esa segunda línea conecta cada pedido con un cliente existente.
\end{frame}

% 20 ---------------------------------------------------------------------
\begin{frame}{\ins{Práctica 3} ¿En qué tabla va la clave foránea?}
\begin{center}
\begin{tikzpicture}
  % caso a
  \node[cajaverde,text width=1.7cm,font=\tiny,align=center] (a1) {\textbf{categoria}};
  \node[cajanaranja,text width=1.7cm,font=\tiny,align=center,right=15mm of a1] (a2) {\textbf{producto}};
  \draw[relacion] (a1)--(a2);
  \pic at ($(a1.east)!2.6mm!(a2.west)$) {uno};
  \pic at ($(a2.west)!2.6mm!(a1.east)$) {muchos};
  % caso b
  \node[cajaverde,text width=1.7cm,font=\tiny,align=center,below=13mm of a1] (b1) {\textbf{autor}};
  \node[cajanaranja,text width=1.7cm,font=\tiny,align=center,right=15mm of b1] (b2) {\textbf{libro}};
  \draw[relacion] (b1)--(b2);
  \pic at ($(b1.east)!2.6mm!(b2.west)$) {uno};
  \pic at ($(b2.west)!2.6mm!(b1.east)$) {muchos};
  % caso c
  \node[cajaverde,text width=1.7cm,font=\tiny,align=center,below=13mm of b1] (c1) {\textbf{area}};
  \node[cajanaranja,text width=1.7cm,font=\tiny,align=center,right=15mm of c1] (c2) {\textbf{empleado}};
  \draw[relacion] (c1)--(c2);
  \pic at ($(c1.east)!2.6mm!(c2.west)$) {uno};
  \pic at ($(c2.west)!2.6mm!(c1.east)$) {muchos};
  % respuestas
  \node[anchor=west,font=\scriptsize,right=14mm of a2] {la columna \bl[1.8cm] va en la tabla \bl[1.8cm]};
  \node[anchor=west,font=\scriptsize,right=14mm of b2] {la columna \bl[1.8cm] va en la tabla \bl[1.8cm]};
  \node[anchor=west,font=\scriptsize,right=14mm of c2] {la columna \bl[1.8cm] va en la tabla \bl[1.8cm]};
\end{tikzpicture}
\end{center}
\vspace{0.4cm}
\small \textbf{Pista:} la clave foránea siempre va del lado de la \textbf{pata de gallo}, y se llama como la otra tabla más \texttt{\_id}.
\end{frame}

% 21 ---------------------------------------------------------------------
\begin{frame}{La tienda de hoy: cuatro tablas, dos conexiones}
\begin{center}
\begin{tikzpicture}
  \node[cajaverde,text width=2.1cm,font=\scriptsize,align=left] (cli)
    {\textbf{cliente}\\ \texttt{id} \pk\\ \texttt{nombre}};
  \node[cajanaranja,text width=2.3cm,font=\scriptsize,align=left,right=24mm of cli] (ped)
    {\textbf{pedido}\\ \texttt{id} \pk\\ \texttt{fecha}\\ \texttt{cliente\_id} \fk};
  \node[cajaverde,text width=2.1cm,font=\scriptsize,align=left,below=12mm of cli] (cat)
    {\textbf{categoria}\\ \texttt{id} \pk\\ \texttt{nombre}};
  \node[cajanaranja,text width=2.3cm,font=\scriptsize,align=left,right=24mm of cat] (pro)
    {\textbf{producto}\\ \texttt{id} \pk\\ \texttt{precio}\\ \texttt{categoria\_id} \fk};
  \draw[relacion] (cli)--(ped); \draw[relacion] (cat)--(pro);
  \pic at ($(cli.east)!2.6mm!(ped.west)$) {uno};
  \pic at ($(ped.west)!2.6mm!(cli.east)$) {muchos};
  \pic at ($(cat.east)!2.6mm!(pro.west)$) {uno};
  \pic at ($(pro.west)!2.6mm!(cat.east)$) {muchos};
  \node[anchor=south,font=\tiny] at ($(cli.east)!.5!(ped.west)$) {hace};
  \node[anchor=south,font=\tiny] at ($(cat.east)!.5!(pro.west)$) {agrupa};
\end{tikzpicture}
\end{center}
\vspace{0.35cm}
\begin{alertblock}{Una tabla no puede apuntar a otra que todavía no existe}
\small Por eso el orden importa: primero las que \textbf{no apuntan a nadie}. \textbf{Numera del 1 al 4}:\\[7pt]
\hspace*{0.5cm}\tbl{producto} \bl[0.9cm] \hspace{0.9cm} \tbl{cliente} \bl[0.9cm] \hspace{0.9cm} \tbl{pedido} \bl[0.9cm] \hspace{0.9cm} \tbl{categoria} \bl[0.9cm]
\end{alertblock}
\end{frame}

% 22 ---------------------------------------------------------------------
\begin{frame}{La tienda completa, de principio a fin en MySQL}
\begin{sqlcode}[.97\textwidth]
\scriptsize
\kw{CREATE DATABASE} tienda; \ \ \kw{USE} tienda;
\end{sqlcode}
\vspace{-0.25cm}
\begin{columns}[T,onlytextwidth]
\column{.49\textwidth}
\begin{sqlcode}[\textwidth]
\scriptsize
\kw{CREATE TABLE} cliente (\\
\ind id INT \kw{PRIMARY KEY} \kw{AUTO\_INCREMENT},\\
\ind nombre VARCHAR(80) \kw{NOT NULL},\\
\ind correo VARCHAR(120) \kw{UNIQUE}\\
);\\[7pt]
\kw{CREATE TABLE} categoria (\\
\ind id INT \kw{PRIMARY KEY} \kw{AUTO\_INCREMENT},\\
\ind nombre VARCHAR(80) \kw{NOT NULL}\\
);
\end{sqlcode}
\column{.49\textwidth}
\begin{sqlcode}[\textwidth]
\scriptsize
\kw{CREATE TABLE} pedido (\\
\ind id INT \kw{PRIMARY KEY} \kw{AUTO\_INCREMENT},\\
\ind fecha DATE \kw{NOT NULL},\\
\ind cliente\_id INT \kw{NOT NULL},\\
\ind \kw{FOREIGN KEY} (cliente\_id)\\
\ind\ind \kw{REFERENCES} cliente(id)\\
);\\[7pt]
\kw{CREATE TABLE} producto (\\
\ind id INT \kw{PRIMARY KEY} \kw{AUTO\_INCREMENT},\\
\ind precio DECIMAL(10,2) \kw{NOT NULL},\\
\ind categoria\_id INT \kw{NOT NULL},\\
\ind \kw{FOREIGN KEY} (categoria\_id)\\
\ind\ind \kw{REFERENCES} categoria(id)\\
);
\end{sqlcode}
\end{columns}
\vspace{0.15cm}
\small \textbf{Todo lo de hoy cabe aquí:} una base, cuatro tablas, sus tipos, sus claves primarias y sus dos claves foráneas.
\end{frame}

% ---------------------------------------------------------------------
\begin{frame}{\ins{Práctica 4} Ejecución del script en MySQL}
\small Copia el código completo de la diapositiva anterior en una nueva pestaña de MySQL Workbench y ejecútalo \textbf{de arriba abajo}.
\vspace{0.45cm}

\begin{sqlcode}[.66\textwidth]
\nuevo{\kw{SHOW TABLES}};\\[7pt]
\nuevo{\kw{DESCRIBE}} cliente;\\
\nuevo{\kw{DESCRIBE}} pedido;
\end{sqlcode}
\vspace{0.35cm}

\begin{block}{Comprueba el resultado}
\small Deben aparecer las tablas \tbl{cliente}, \tbl{categoria}, \tbl{pedido} y \tbl{producto}. Revisa que \texttt{id} sea la clave primaria y que \texttt{cliente\_id} aparezca en \tbl{pedido}.
\end{block}
\vspace{0.3cm}

\small Si aparece un error, identifica la instrucción que falló, corrígela y vuelve a ejecutar esa instrucción.
\end{frame}

% ---------------------------------------------------------------------
\begin{frame}{La misma tienda, en \textnormal{SQL Server}}
\begin{sqlcode}[.97\textwidth]
\scriptsize
\kw{CREATE DATABASE} tienda; \ \ \nuevo{GO} \ \ \kw{USE} tienda; \ \ \nuevo{GO}
\end{sqlcode}
\vspace{-0.25cm}
\begin{columns}[T,onlytextwidth]
\column{.49\textwidth}
\begin{sqlcode}[\textwidth]
\scriptsize
\kw{CREATE TABLE} cliente (\\
\ind id INT \nuevo{\kw{IDENTITY}(1,1)} \kw{PRIMARY KEY},\\
\ind nombre \nuevo{NVARCHAR(80)} \kw{NOT NULL},\\
\ind correo \nuevo{NVARCHAR(120)} \kw{UNIQUE}\\
);\\[7pt]
\kw{CREATE TABLE} categoria (\\
\ind id INT \nuevo{\kw{IDENTITY}(1,1)} \kw{PRIMARY KEY},\\
\ind nombre \nuevo{NVARCHAR(80)} \kw{NOT NULL}\\
);
\end{sqlcode}
\column{.49\textwidth}
\begin{sqlcode}[\textwidth]
\scriptsize
\kw{CREATE TABLE} pedido (\\
\ind id INT \nuevo{\kw{IDENTITY}(1,1)} \kw{PRIMARY KEY},\\
\ind fecha DATE \kw{NOT NULL},\\
\ind cliente\_id INT \kw{NOT NULL},\\
\ind \kw{FOREIGN KEY} (cliente\_id)\\
\ind\ind \kw{REFERENCES} cliente(id)\\
);\\[7pt]
\kw{CREATE TABLE} producto (\\
\ind id INT \nuevo{\kw{IDENTITY}(1,1)} \kw{PRIMARY KEY},\\
\ind precio DECIMAL(10,2) \kw{NOT NULL},\\
\ind categoria\_id INT \kw{NOT NULL},\\
\ind \kw{FOREIGN KEY} (categoria\_id)\\
\ind\ind \kw{REFERENCES} categoria(id)\\
);
\end{sqlcode}
\end{columns}
\vspace{0.15cm}
\small \textbf{Lo marcado es lo único que cambia:} \texttt{IDENTITY(1,1)} en lugar de \texttt{AUTO\_INCREMENT}, \texttt{NVARCHAR} en lugar de \texttt{VARCHAR}, y \texttt{GO} para separar instrucciones. \textbf{Todo lo demás es idéntico.}
\end{frame}

% ---------------------------------------------------------------------
\begin{frame}{La misma tienda, en \textnormal{SQLite}}
\begin{sqlcode}[.97\textwidth]
\scriptsize
\nuevo{-- no hay CREATE DATABASE ni USE: la base de datos \textbf{es un archivo}} \ \ \nuevo{\kw{PRAGMA} foreign\_keys = ON;}
\end{sqlcode}
\vspace{-0.25cm}
\begin{columns}[T,onlytextwidth]
\column{.49\textwidth}
\begin{sqlcode}[\textwidth]
\scriptsize
\kw{CREATE TABLE} cliente (\\
\ind id \nuevo{INTEGER} \kw{PRIMARY KEY} \nuevo{\kw{AUTOINCREMENT},}\\
\ind nombre \nuevo{TEXT} \kw{NOT NULL},\\
\ind correo \nuevo{TEXT} \kw{UNIQUE}\\
);\\[7pt]
\kw{CREATE TABLE} categoria (\\
\ind id \nuevo{INTEGER} \kw{PRIMARY KEY} \nuevo{\kw{AUTOINCREMENT},}\\
\ind nombre \nuevo{TEXT} \kw{NOT NULL}\\
);
\end{sqlcode}
\column{.49\textwidth}
\begin{sqlcode}[\textwidth]
\scriptsize
\kw{CREATE TABLE} pedido (\\
\ind id \nuevo{INTEGER} \kw{PRIMARY KEY} \nuevo{\kw{AUTOINCREMENT},}\\
\ind fecha \nuevo{TEXT} \kw{NOT NULL},\\
\ind cliente\_id INTEGER \kw{NOT NULL},\\
\ind \kw{FOREIGN KEY} (cliente\_id)\\
\ind\ind \kw{REFERENCES} cliente(id)\\
);\\[7pt]
\kw{CREATE TABLE} producto (\\
\ind id \nuevo{INTEGER} \kw{PRIMARY KEY} \nuevo{\kw{AUTOINCREMENT},}\\
\ind precio \nuevo{REAL} \kw{NOT NULL},\\
\ind categoria\_id INTEGER \kw{NOT NULL},\\
\ind \kw{FOREIGN KEY} (categoria\_id)\\
\ind\ind \kw{REFERENCES} categoria(id)\\
);
\end{sqlcode}
\end{columns}
\vspace{0.15cm}
\small SQLite tiene \textbf{muy pocos tipos}: \texttt{INTEGER}, \texttt{TEXT} y \texttt{REAL} cubren todo, y las fechas se guardan como texto. \textbf{Y las claves foráneas vienen apagadas}: sin el \texttt{PRAGMA} se aceptan y no se vigilan.
\end{frame}

% ---------------------------------------------------------------------
\begin{frame}{Cada gestor tiene su propia notación}
\begin{block}{Lo que hay que llevarse}
El \textbf{SQL que aprendimos hoy es el mismo en todos}. Lo que cambia son \textbf{detalles}: cómo se numera solo el \texttt{id}, cómo se llama el tipo de texto, si existe o no \texttt{CREATE DATABASE}.
\end{block}
\vspace{0.4cm}
\begin{center}
\small
\begin{tabular}{>{\raggedright\arraybackslash}p{3.1cm}>{\ttfamily}p{3.0cm}>{\ttfamily}p{2.9cm}>{\ttfamily}p{2.4cm}}
\toprule
& \normalfont\bfseries MySQL & \normalfont\bfseries SQL Server & \normalfont\bfseries SQLite \\
\midrule
Numerar solo el \texttt{id} & AUTO\_INCREMENT & IDENTITY(1,1) & AUTOINCREMENT \\[4pt]
Texto & VARCHAR(80) & NVARCHAR(80) & TEXT \\[4pt]
Decimales & DECIMAL(10,2) & DECIMAL(10,2) & REAL \\[4pt]
Crear la base & CREATE DATABASE & CREATE DATABASE & \normalfont es un archivo \\
\bottomrule
\end{tabular}
\end{center}
\vspace{0.4cm}
\small Por eso no hay que aprender tres idiomas: se aprende \textbf{uno}, y al cambiar de gestor se consulta la tabla de equivalencias. \texttt{CREATE TABLE}, \texttt{PRIMARY KEY}, \texttt{FOREIGN KEY}, \texttt{NOT NULL} y \texttt{UNIQUE} son iguales en los tres.
\end{frame}

% ---------------------------------------------------------------------
\begin{frame}{\ins{Práctica 5} Ahora prueba guardar y consultar}
\small Las tablas ya existen, pero están \textbf{vacías}. Escribe estas dos instrucciones y mira qué responde el gestor:
\vspace{0.3cm}
\begin{sqlcode}[.80\textwidth]
\textcolor{utpGris}{-- 1. guardar una fila}\\
\kw{INSERT INTO} cliente (nombre, correo)\\
\kw{VALUES} ('Raúl', 'raul@mail.com');\\[9pt]
\textcolor{utpGris}{-- 2. verla}\\
\kw{SELECT} * \kw{FROM} cliente;
\end{sqlcode}
\vspace{0.15cm}
\small \texttt{INSERT INTO} dice \textbf{en qué tabla y en qué columnas}; \texttt{VALUES} trae los datos, en el mismo orden. \texttt{SELECT * FROM} es «muéstrame todo de esta tabla».
\vspace{0.3cm}
\begin{block}{Dos cosas para comprobar}
\small ¿Qué valor le tocó al \texttt{id}, si no lo escribimos? \ Y prueba un pedido con \texttt{cliente\_id = 97}: \textbf{¿qué contesta?}
\end{block}
\end{frame}

% ---------------------------------------------------------------------
\begin{frame}{Lo de hoy fue DDL: ¿y el resto?}
\begin{center}
\resizebox{\textwidth}{!}{%
\begin{tikzpicture}
  \node[cajanaranja,text width=3.1cm,font=\scriptsize,align=center] (ddl)
    {\textbf{DDL}\\[2pt] \textit{definición}\\[2pt] define la estructura\\[3pt] \texttt{CREATE ALTER DROP}};
  \node[caja,text width=3.1cm,font=\scriptsize,align=center,right=9mm of ddl] (dml)
    {\textbf{DML}\\[2pt] \textit{manipulación}\\[2pt] cambia los datos\\[3pt] \texttt{INSERT UPDATE DELETE}};
  \node[caja,text width=3.1cm,font=\scriptsize,align=center,right=9mm of dml] (dql)
    {\textbf{DQL}\\[2pt] \textit{consulta}\\[2pt] lee los datos\\[3pt] \texttt{SELECT}};
  \node[caja,text width=3.1cm,font=\scriptsize,align=center,right=9mm of dql] (dcl)
    {\textbf{DCL}\\[2pt] \textit{control}\\[2pt] quién puede qué\\[3pt] \texttt{GRANT REVOKE}};
  \node[font=\scriptsize,below=3mm of ddl,text=utpRojo] {\textbf{todo lo de hoy}};
  \node[font=\scriptsize,below=3mm of dml] {la práctica 4};
  \node[font=\scriptsize,below=3mm of dql] {próximas sesiones};
  \node[font=\scriptsize,below=3mm of dcl] {más adelante};
\end{tikzpicture}}
\end{center}
\vspace{0.5cm}
\small Las cuatro son SQL; se agrupan por \textbf{para qué sirven}. Hoy solo construimos el \textbf{envase}: qué tablas hay y qué columnas tienen. Con el envase hecho, lo que sigue es \textbf{llenarlo y consultarlo}.
\end{frame}

% 23 ---------------------------------------------------------------------
\begin{frame}{Cierre: completa en voz alta}
\begin{itemize}\setlength\itemsep{16pt}
  \item Una columna se escribe siempre con dos cosas: \bl[5cm]
\end{itemize}
\end{frame}

% 24 ---------------------------------------------------------------------
\begin{frame}{Referencias}
\footnotesize
Coronel, C., Morris, S., \& Rob, P. (2011). \textit{Bases de datos: diseño, implementación y administración} (9.ª ed.). Cengage Learning.

\vspace{0.35cm}
Elmasri, R., \& Navathe, S. (2007). \textit{Fundamentos de sistemas de bases de datos} (5.ª ed.). Addison Wesley.

\vspace{0.35cm}
Silberschatz, A., Korth, H. F., \& Sudarshan, S. (2020). \textit{Database system concepts} (7th ed.). McGraw-Hill.
\end{frame}

% ---------------------------------------------------------------------
% Cierre UTP
{
  \setbeamertemplate{background canvas}{%
    \includegraphics[width=\paperwidth,height=\paperheight]{gracias-final.pdf}%
  }
  \setbeamertemplate{footline}{}
  \begin{frame}[plain]
  \end{frame}
}

\end{document}
